\subsection{Local corrections need exact recourse}\label{sec:local}
The correction of Section~\ref{sec:grids} is global: every coordinate is
rounded, and the allowance is the sum of all coordinate corrections. A
bag-local allowance is valid when the coordinates outside the bag are
optimized exactly or with a certified error.

\begin{theorem}[Conditional-recourse filter]\label{thm:cr-filter}
Let $\mathcal Q$ be a finite family of bag cells and
$e\ge\max_{C\in\mathcal Q}e_B(C)$. Suppose that for every corner $v$ of
every cell of $\mathcal Q$ an oracle supplies a number $\ell(v)$ and a
point $x^v\in X$ with $x^v_B=v$ such that
\begin{equation}\label{eq:cr-oracle}
\ell(v)\le W_B(v),\qquad F(x^v)\le\ell(v)+\eta .
\end{equation}
Let $U$ be the value of a feasible point with $U\le\min_vF(x^v)$, the
minimum taken over all supplied completions. Retain $C\in\mathcal Q$ if
and only if $\min_{v\text{ corner of }C}\ell(v)-e\le U$. Then:
\begin{enumerate}
\item[(a)] every $C\in\mathcal Q$ that contains $x^*_B$ for some global
optimizer $x^*$ is retained;
\item[(b)] if some $C\in\mathcal Q$ contains $x^*_B$ for a global optimizer
$x^*$, then $U\le\OPT+e+\eta$, and every retained cell has a corner $v$
with $W_B(v)\le\OPT+2e+2\eta$;
\item[(c)] every $x\in X$ with $x_B\in\bigcup\mathcal Q$ satisfies
\[
F(x)\ge\lambda:=\min_{C\in\mathcal Q}\Bigl(\min_{\text{corners of }C}\ell-e\Bigr),
\]
and $U-\lambda\le e+\eta$.
\end{enumerate}
\end{theorem}
\begin{proof}
(a) By~\eqref{eq:cr-oracle} and Lemma~\ref{lem:cr-cell} applied to $x^*$,
\[
\min_{\text{corners}}\ell-e\le\min_{\text{corners}}W_B-e_B(C)\le\OPT\le U.
\]

(b) Lemma~\ref{lem:cr-cell} at $x^*$ gives a corner $v^\circ$ of the
optimizer's cell with $W_B(v^\circ)\le\OPT+e$. Hence
$U\le F(x^{v^\circ})\le\ell(v^\circ)+\eta\le W_B(v^\circ)+\eta\le\OPT+e+\eta$.
A retained cell has a corner $v$ with $\ell(v)\le U+e$, and
$W_B(v)\le F(x^v)\le\ell(v)+\eta\le U+e+\eta\le\OPT+2e+2\eta$.

(c) For $x_B\in C$, Lemma~\ref{lem:cr-cell} and $\ell\le W_B$ give
$F(x)\ge\min_{\text{corners}}W_B-e\ge\min_{\text{corners}}\ell-e\ge\lambda$.
Every corner satisfies $\ell(v)\ge F(x^v)-\eta\ge U-\eta$, so
$\lambda\ge U-e-\eta$.
\end{proof}

\begin{remark}
Hypothesis~\eqref{eq:cr-oracle} is a certified global optimization of the
outside problem, over the original outside domain, at every queried bag
value, with an error $\eta$ that is uniform over these values. In a
multilevel scheme that refines only retained cells, the hypothesis of
(b) holds at every level by (a) and induction. A point whose bag
projection lies in a cell discarded at an earlier level has value larger
than the incumbent at that level, hence larger than the current $U$.
Therefore $\min\{\lambda,U\}$ is a certified lower bound on $\OPT$, as in
the filtering proposition for corrected grids.
\end{remark}

The next proposition shows that the filter needs only control of the
\emph{variation} of the error across bag states. A common additive error,
such as a normalization constant of messages, is harmless.

\begin{proposition}[Only the error oscillation matters]\label{prop:cr-osc}
Let $\mathcal Q$, $e$ be as in Theorem~\ref{thm:cr-filter}, and suppose
some cell of $\mathcal Q$ contains $x^*_B$ for a global optimizer $x^*$.
Let $\mathcal V$ be the set of all corners and $\ell:\mathcal V\to\R$ a
function such that, for some constants $\delta_-\le\delta_+$ with
$\delta_+-\delta_-\le\eta$,
\[
\delta_-\le W_B(v)-\ell(v)\le\delta_+\qquad(v\in\mathcal V).
\]
The constants need not be known. Retain $C$ if and only if
$\min_{\text{corners of }C}\ell\le\min_{\mathcal V}\ell+e+\eta$. Then every
cell containing $x^*_B$ for a global optimizer $x^*$ is retained, and
every retained cell has a corner with $W_B(v)\le\OPT+2e+2\eta$.
\end{proposition}
\begin{proof}
Let $C^*$ contain $x^*_B$. By Lemma~\ref{lem:cr-cell},
\[
\min_{\text{corners of }C^*}\ell\le\min_{\text{corners of }C^*}W_B-\delta_-
\le\OPT+e-\delta_-,
\]
while $W_B\ge\OPT$ gives
$\min_{\mathcal V}\ell\ge\OPT-\delta_+$. The difference is at most
$e+\delta_+-\delta_-\le e+\eta$, so $C^*$ is retained. If $C$ is retained,
it has a corner $v$ with
$W_B(v)\le\ell(v)+\delta_+\le\min_{\mathcal V}\ell+e+\eta+\delta_+$, and
$\min_{\mathcal V}\ell\le\min_{\text{corners of }C^*}\ell\le\OPT+e-\delta_-$.
Thus $W_B(v)\le\OPT+2e+\eta+(\delta_+-\delta_-)\le\OPT+2e+2\eta$.
\end{proof}

Theorem~\ref{thm:cr-filter} additionally certifies a gap; for that, valid
lower bounds and feasible completions are needed. For filtering alone,
Proposition~\ref{prop:cr-osc} identifies the essential requirement: the
oscillation of $W_B-\ell$ over the compared bag states must be
$O(e)$.

\paragraph{Grid min-marginals are not certified recourse.}
Ordinary grid min-marginals optimize the outside coordinates on the grid,
not over their original domain. The next proposition shows that no
allowance depending only on the bag, the curvature and the mesh can make
them valid.

Let $h=2^{-j}$ and $G=h\Z\cap[0,1]$. On $[0,1]^{1+m}$ with coordinates
$(x,y_1,\ldots,y_m)$, let the uniform product grid be $G^{1+m}$, write
$U_G=\min_{G^{1+m}}F$, and, for the bag $B=\{x,y_1\}$ and $v\in G^2$,
\[
m^F(v)=\min\{F(z):z\in G^{1+m},\ (z_x,z_{y_1})=v\}.
\]
A \emph{constant-correction test} with constant $c$ retains a bag cell
$C$ of the grid if and only if $\min_{\text{corners of }C}m^F-c\le U_G$;
its coordinate version retains a grid interval $[a,a+h]$ of the centre
coordinate if and only if $\min\{V_h(a),V_h(a+h)\}-c\le U_G$, where
$V_h(a)=\min\{F(z):z\in G^{1+m},\ z_x=a\}$. On a uniform grid every node has
correction $Lh^2/8$ in every coordinate, so the corrected-grid method of
the paper applies the coordinate version with $c=(1+m)Lh^2/8$. A bag-local
budget would take $c=|B|Lh^2/8=Lh^2/4$. In both families below, the
corner minimum of each relevant bag cell equals
$\min\{V_h(a),V_h(a+h)\}$, where $[a,a+h]$ is the projection of the cell
on the centre coordinate. Hence the thresholds stated for bag cells hold
verbatim for the coordinate version.

\begin{proposition}[Constant corrections must grow with dimension]\label{prop:star}
\leavevmode
\begin{enumerate}
\item[(A)] Let $j\ge1$, $m\ge1$, $\epsilon>0$ with $mh^2>4(1+\epsilon)$, and
\[
F_A(x,y)=x^2-hx\sum_{i=1}^my_i+\sum_{i=1}^my_i^2
+\Bigl(\frac{mh^2}4-1-\epsilon\Bigr)x .
\]
Its interaction graph is a star with bags $\{x,y_i\}$ ($p=2$), and $L=2$
is valid. Its unique minimizer is $x=1$, $y_i=h/2$, with $\OPT=-\epsilon$,
and any growth constant satisfies $g\le\epsilon/(1+mh^2/4)$; moreover
$U_G=0$. The only bag cell containing $(1,h/2)$ is
$C=[1-h,1]\times[0,h]$, and
\[
\min_{\text{corners of }C}m^F=(1-h)\Bigl(\frac{mh^2}4-h-\epsilon\Bigr).
\]
Hence a constant-correction test discards a cell containing the unique
optimizer unless
$c\ge(1-h)\bigl(m-4/h-4\epsilon/h^2\bigr)\,Lh^2/8$. As $m\to\infty$ the
ratio of $(1+m)Lh^2/8$ to this lower bound tends to $(1-h)^{-1}$.
\item[(B)] Let $j\ge2$, let $d\ge2/h$ be an integer, $m=d^2$, and
\[
F_B(x,y)=\Bigl(x-\frac12\Bigr)^2+\sum_{i=1}^m
\Bigl(y_i-\frac h2-\frac{x-1/2}{d}\Bigr)^2 .
\]
It is a star with $p=2$; $L=4$ is valid; its unique minimizer
$z^*=(1/2,h/2,\ldots,h/2)$ is interior; and
$F_B(z)-\OPT\ge\frac{3-\sqrt5}2\norm{z-z^*}^2$, so
$\kappa\le2(3+\sqrt5)<11$ for every $d$ and $h$. Moreover
$U_G=\frac12-\frac{dh}2+\frac{d^2h^2}4$, and both bag cells containing
$(1/2,h/2)$ have
\[
\min_{\text{corners}}m^F-U_G=dh\Bigl(\frac12-h\Bigr)+2h^2-\frac12 .
\]
Thus a constant-correction test discards every cell containing the
unique optimizer unless $c\ge dh(\frac12-h)+2h^2-\frac12$, a quantity of
order $\sqrt m\,h$ at fixed $h$, while $p=2$ and $\kappa<11$.
\end{enumerate}
\end{proposition}
\begin{proof}
(A) Write $\beta=mh^2/4-1-\epsilon>0$. All coordinate second derivatives
equal $2$ and $F_A$ is quadratic, so $L=2$ is valid. For fixed $x$, each
leaf term $y^2-hxy$ is strictly convex with minimizer $hx/2\in[0,1]$ and
minimum $-h^2x^2/4$. Hence $\min_yF_A(x,y)=(1-mh^2/4)x^2+\beta x=:V(x)$,
which is strictly concave because $mh^2>4$, with $V(0)=0$ and
$V(1)=-\epsilon$. For $x\in(0,1)$, strict concavity gives
$V(x)>(1-x)V(0)+xV(1)\ge-\epsilon$. So $x=1$ is the unique minimizing
centre and $y_i=h/2$ the unique leaves; $\OPT=-\epsilon$. The origin has
$F_A=0=\OPT+\epsilon$ and squared distance $1+mh^2/4$ from the minimizer,
which bounds $g$.

On the grid, $y(y-hx)\ge0$ for $y\in G$ and $x\in[0,1]$, since either
$y=0$ or $y\ge h\ge hx$. Also $x^2+\beta x\ge0$ for $x\ge0$. Hence
$F_A\ge0$ on $G^{1+m}$ with equality at the origin, so $U_G=0$, and every
leaf's grid minimum is $0$. The point $(1,h/2)$ has $x$ at the right
endpoint and $h/2$ interior to $[0,h]$, so $C$ is the only grid cell of
the bag containing it. At its corners,
$m^F(x,y_1)=x^2+\beta x+(y_1^2-hxy_1)$. For $x=1-h$ this equals
$(1-h)(1-h+\beta)=(1-h)(mh^2/4-h-\epsilon)$ at $y_1=0$ and exceeds it by
$h^3$ at $y_1=h$. For $x=1$ it equals $mh^2/4-\epsilon$ at both $y_1=0,h$,
which exceeds the previous value by $h(mh^2/4-\epsilon)+h(1-h)>0$. This
proves the corner minimum; it also equals $\min\{V_h(1-h),V_h(1)\}$,
because every leaf's grid minimum is attained at $0$. Since $U_G=0$ and
$L=2$, the threshold for $c$ is as stated, and
$(1+m)h^2/4\big/\bigl[(1-h)(mh^2/4-h-\epsilon)\bigr]\to(1-h)^{-1}$.

(B) The Hessian $H$ has $H_{xx}=2+2m/d^2=4$, $H_{y_iy_i}=2$,
$H_{xy_i}=-2/d$, and zero elsewhere; thus $L=4$ is valid. Vectors with
zero $x$-component and leaf components summing to zero are eigenvectors
with eigenvalue $2$. On the span of $e_x$ and $m^{-1/2}\mathbf 1_y$, $H$
acts as $\bigl(\begin{smallmatrix}4&-2\\-2&2\end{smallmatrix}\bigr)$, with
eigenvalues $3\pm\sqrt5$. Since $F_B\ge0$ vanishes only at $z^*\in(0,1)^{1+m}$
and $F_B(z)=\frac12(z-z^*)^{\T}H(z-z^*)$, growth holds with
$g=(3-\sqrt5)/2$ and $\kappa=L/g=2(3+\sqrt5)$.

Write $t=x-1/2\in h\Z\cap[-1/2,1/2]$ for centre grid values (here
$1/2\in G$ because $j\ge1$). A leaf's ideal value
$\iota(t)=h/2+t/d$ lies in $[0,h]$ because $|t|/d\le1/(2d)\le h/4$.
Its distance to $G$ is therefore $h/2-|t|/d$, and the conditional grid
value of the centre is
\[
V_h(t)=t^2+m\Bigl(\frac h2-\frac{|t|}d\Bigr)^2=2t^2-dh|t|+\frac{d^2h^2}4 .
\]
On the grid $|t|\le1/2\le dh/4$, and $s\mapsto2s^2-dhs$ decreases on
$[0,dh/4]$, so $U_G=V_h(\pm1/2)=\frac12-\frac{dh}2+\frac{d^2h^2}4$.
The optimizer's bag projection $(1/2,h/2)$ lies on the common edge of
the cells $[1/2-h,1/2]\times[0,h]$ and $[1/2,1/2+h]\times[0,h]$ and in no
other cell. At $t=0$ both leaf labels $0,h$ are at distance $h/2$ from
$\iota$, so $m^F=d^2h^2/4$. At $t=\pm h$ the nearer label gives
$m^F=V_h(h)=2h^2-dh^2+d^2h^2/4$ and the farther adds $2h^2/d$. Since
$d\ge2$, each cell has corner minimum $d^2h^2/4-(d-2)h^2
=\min\{V_h(0),V_h(\pm h)\}$, and subtracting $U_G$ gives the stated
difference. For $h\le1/4$ this difference is at least
$dh/4+2h^2-1/2$, which grows linearly in $d=\sqrt m$. Discarding both
centre intervals adjacent to $x=1/2$ removes the optimizer.
\end{proof}

\begin{example}[The smallest instance]\label{ex:cr-star32}
Take $h=1/2$, $m=32$, $\epsilon=1/16$ in (A):
$F=x^2-\frac x2\sum y_i+\sum y_i^2+\frac{15}{16}x$ on $[0,1]^{33}$.
The exact leaf recourse is $y_i=x/4$, with $V(x)=-x^2+\frac{15}{16}x$, so
the unique optimizer is $x=1$, $y_i=1/4$, value $-1/16$. On
$G=\{0,\frac12,1\}$ the conditional grid values at $x=0,\frac12,1$ are
$0,\frac{23}{32},\frac{31}{16}$, and $U_G=0$. The four corner
min-marginals of $C=[\frac12,1]\times[0,\frac12]$ are
$\frac{23}{32},\frac{27}{32},\frac{31}{16},\frac{31}{16}$. The bag-local
budget $\frac18$ rejects $C$; the global correction $\frac{33}{16}$ keeps it.
The conditional grid error is
\[
V_h(x)-V(x)=m\,\dist(x/4,G)^2,
\]
which varies from $0$ at $x=0$ to $m/16$ at $x=1$; no normalization of
messages removes this variation. Because the growth constant of family (A)
degrades as $m$ grows, this instance refutes budgets depending on
$(p,L)$ only; family (B) refutes budgets depending on $(p,\kappa)$.
\end{example}

